\summarychapter{29}{Opening the Rock}{America and Asia}

\begin{inbrief}
Between 1500 and 1650, while Europe split into Catholic and Protestant, Roman Catholic monks and
friars carried the faith to the Spanish colonies of Latin America and the Portuguese ports of Africa
and Asia. The chapter follows this ``age of global expansion'' through the conquistadores and Las
Casas, Francis Xavier in India and Japan, and Matthew Ricci in China, and the lasting tension between
two missionary policies, conquest and adaptation.
\question{Should the peoples of the new lands be won to the faith by love or by force? How does the
gospel relate to culture?}
\end{inbrief}

On Thursday, October 11, 1492, Christopher Columbus sighted land in what we call the West Indies. The
next morning he named the island San Salvador, ``Holy Savior,'' claimed it for Ferdinand and Isabella
of Spain and gave the people red caps and glass beads, judging from their pleasure that they could be
``converted to our holy faith by love'' rather than by force. That alternative troubled missionaries
through the whole age of discovery. Sometimes a selfless Franciscan or Jesuit won people by patient
love; sometimes conquerors forced baptism at sword point. Behind it lay the question of what
heathenism is: humanity stumbling in search of the true God, or humanity organized in resistance to
the gospel? Should the missionary look for the good in other religions and build on it, the policy of
adaptation, or suppress them in order to plant the true faith, the policy of conquest?

\ssection{The Gospel and Culture}
The church is always in danger of identifying the gospel with the cultural form in which it has made
its home. The first Jewish Christians believed the faith inseparable from Hebrew rites and sabbaths,
and only when it broke free could it take root in the Greco-Roman world; it had to find new ways again
when the Germanic tribes destroyed Roman institutions. In the Middle Ages the rise of Islam cut
European Christians off from their fellow believers in Asia and Africa and confined Christianity
almost entirely to Europe. The crusades were Europe's convulsive attempt to break through by force,
and the long war to drive the Moors from Spain and Portugal bred an especially militant attitude
toward the unbeliever. Conquest and evangelization became intertwined, and few saw a difference.

The new age of missions opened with the great discoveries. From the middle of the fifteenth century
the Portuguese and Spaniards sailed south, soon followed by the French, Dutch and English. The
Portuguese reached the Cape of Good Hope in 1486, and Vasco da Gama went on around Africa to the
Malabar Coast of India, where Arab merchants had for centuries collected spices, gems, ivory, cottons
and silks for the Mediterranean. The Portuguese drove the Arab ships from the Indian Ocean, made Goa
the center of their eastern empire, and pushed on to the Malay Archipelago and the Moluccas, reaching
China by 1516 and Japan by 1543. Unlike the Spanish they sought no large conquests or colonies. They
wanted a monopoly of trade and therefore control of the seas. Their trading ports, each with a small Catholic church, dotted the coasts of West Africa, the Congo, Angola and Mozambique, of India, Ceylon and Malaya, and of Brazil.

\ssection{Spaniards in the New World}
Columbus made three more voyages in a vain search for a way to Asia. On the second, with seventeen
ships and 1,500 men, he found gold on Haiti, and despite an oath of secrecy the news set off the most
violent gold rush in history. Half a century was enough for the Spaniards to overrun and loot the Americas from California to the far south of the continent; the word \emph{conquistadores} comes from these
brave and brutal men. By 1521 Hernando Cortes, with horses, armor and gunpowder, had destroyed the
Aztecs' empire in Mexico; by 1533 Francisco Pizarro had, by treachery, killed Atahualpa and brought
the Inca empire to its knees. All this happened in Luther's lifetime. Charles V, before whom Luther
stood at Worms, was heir to this wealth and took comfort in the treasure ships while Germany slipped
from him.

It was not only a race for spices and gold. Merchants and conquerors alike were zealous to extend the
Roman Catholic faith, and hardly a ship left without friars of the Dominican, Franciscan or Augustinian orders, or Jesuits, on board. Their sacrifices were extraordinary. There were no furloughs, and a third of the Jesuits who set out for China between 1581 and 1712 (127 of 376) died on the way. So the century in which ``catholic'' ceased to mean
united saw Christianity become catholic in a new sense, in its mission to distant peoples. To prevent
rivalry the pope drew a line from pole to pole west of the Azores: everything to the west belonged to
Spain, everything to the east to Portugal, which is why Brazil speaks Portuguese. By papal commission
the two kings were to evangelize the lands their soldiers conquered. They sent and maintained the
missionaries, controlled them exclusively and named the bishops. The governing idea was an extension
of the Christian prince, the divinely appointed guide who removes idolatry from his lands and brings
his subjects under Christ through the Catholic Church.

Spain's devotion to the church and its long struggle with Islam gave its expeditions the fanaticism
of a crusade. Cortes carried a statuette of the Virgin, prayed and heard mass daily, and marched
under flags bearing the cross and Mary's picture. Spain's first policy, the \emph{encomienda},
assigned Indians to colonists to work mines and plantations in return for protection and instruction
in the faith. Because the Indians practiced idolatry and human sacrifice, the Spanish thought their
wars righteous, like Israel's against the Canaanites. In 1531 Bishop Zumarraga reported from Mexico
that more than five hundred temples and twenty thousand idols had been destroyed.

\ssection{Champion of the Oppressed}
The Indians' champion was Bartholomew de Las Casas, whose father had sailed with Columbus on the
second voyage. At first he shared in Christian imperialism like any other colonist. In his middle
thirties he was converted and ordained, the first ordination in the New World, and began to preach
against the cruelties around him: how could God ordain brutality as a means of conversion? Many
thought him a visionary and a menace. Some missionaries held that barbarous peoples cannot be brought
to the faith without force, and colonists warned that ending the encomienda would ruin the colonies.
Las Casas crossed the Atlantic fourteen times to plead his cause. Partly through him Charles V issued
the Laws of the Indies in 1542, which softened the system and established the principle that Indians
too have human rights. But law in Spain was one thing and practice in Mexico another. In a famous
debate at Valladolid in 1550 Las Casas argued for the Indians' equality and freedom and for conversion
only by peaceful preaching and holy example. His opponent, the theologian Juan Gin\'es de Sep\'ulveda,
appealed to Aristotle's teaching that some peoples are born to slavery. Christian imperialism went
on and hampered the gospel.

The Indians accepted the faith readily and abandoned it as readily. In seven years, from 1524 to 1531, Franciscan friars baptized over a million people in Mexico, a mass baptism recalling the early Middle Ages.
One story tells how Cortes left a town with orders to worship the Christian God and to care for his
lame horse; the people fed the horse fruit and flowers until it died, and Franciscans later found
them worshiping its image as a god of storms. The Portuguese in the East likewise
held the king responsible for his subjects' souls. Goa became a city of baroque churches with an
archbishop for all East Asia, and no idolatry was allowed on Portuguese territory. But Asia confronted them with advanced civilizations and old religions of a strength the Aztecs and Incas had never had. In India, Japan and China the policy of conquest was modified or abandoned for adaptation, and
the Jesuits in particular asked which local customs were merely social and which were incompatible
with baptism.

\ssection{Francis Xavier}
The pioneer in the Far East, Francis Xavier (1506--1552), Loyola's associate, landed at Goa in 1542
a legate of the pope employed by Portugal's king. In a few months he stirred a moral revolution
in that rough seaport, then went to the southern tip of India. There, in 1534, the Paravars, some
thirty thousand fisherfolk harassed by robbers, had accepted baptism in return for Portuguese
protection, and for eight years no one had taught them. Xavier went from village to village across
the burning sand, ringing a handbell, reciting Creed, Commandments, Lord's Prayer and Rosary, and baptizing by the hundreds until his arms gave out. In three years he left an
organized church. Because it was confined to one caste, it remained an island in the sea of Hinduism,
but the Jesuits who followed gathered the Paravars into villages centered on the church, and in time
all memory of their not being Christian faded.

Catholics and Protestants alike praise Xavier's zeal and charm, but he was a man of his age. In 1546
he asked the king of Portugal to set up the Inquisition in Goa against the Muslims, and it lasted
there until the early nineteenth century. That year he moved on to Malaya, and in 1549 he sailed for
Japan, landing on August 15 with two companions. The moment was favorable. Japan was ruled by feudal
barons eager for trade, Buddhism was out of favor, and no national religion stood in the way. Japan
changed his mind about missions. In India he had held, like most missionaries in Latin America, that
heathenism has nothing in common with the faith and must be leveled before anything Christian can be
built. Among a people with so many marks of nobility he saw that the gospel must transform and refine
but need not discard everything in their life and culture. The Jesuits who followed developed a
quality quite untypical of the Counter Reformation, compromise in religious matters. They adapted to
local customs wherever they could and trained Japanese leaders as fast as possible. In 1579 they
founded Nagasaki as a home for converts, and by the end of the century they counted three hundred
thousand converts, hundreds of churches and two colleges. Early in the seventeenth century new
rulers, fearing that foreign missionaries meant foreign invaders, turned to persecution. Between 1614
and 1642 there were 4,045 martyrs, by sword, stake and boiling cauldron, and only a remnant survived
in the hills near Nagasaki.

\ssection{The Door to China}
After two years in Japan Xavier set out for China. Hoping to smuggle himself into Canton, he reached
an island off the south coast, fell ill and died there in 1552. The door opened to Matthew Ricci
(1552--1610). Macao became a Portuguese colony in 1567, but the Ming dynasty wanted no contact with
the outside world. The Chinese thought of themselves as givers of culture, and Confucianism governed
state, family and ethics. In 1579 the Jesuit leader Alessandro Valignani, looking toward the coast
from his window, cried, ``Oh, Rock, Rock, when will you open?'' He sent for Ricci, an Italian Jesuit
trained in mathematics, astronomy and cosmology. Ricci learned the language and customs at Macao and
in 1583 was allowed to settle in the provincial capital, Chaoch'ing. Dressed as a Mandarin, he won
the respect the Chinese give the scholar, made them a map of the world and introduced the new science
of the calendar. Convinced that China could not be won from a province, he moved step by step to
Peking, which he entered in 1600. His gift of two clocks delighted the emperor, and since only Ricci
could keep them running he remained ten years as astronomer and mathematician. Notable families and
scholars were baptized, and at his death the church numbered two thousand.

His successor Adam Schall won the scholars' admiration by predicting an eclipse of the moon and
became director of the Imperial Astronomical Service. He built a public church in Peking in 1650 and
gained freedom for Christianity throughout the empire in 1657; at his death there were almost 270,000
Christians. An imperial edict of toleration in 1692 rewarded the Jesuits' services. The mission was
then damaged from within. Ricci had presented Christianity not as something new but as the perfection
of traditional Chinese devotion: the ``Lord of Heaven'' long revered was God, and reverence for
ancestors was a social, not a religious, act. A Franciscan and a Dominican who came to Peking in 1631
were shocked to find the mass translated by the word for the ancestor ceremony and, after attending
such a ceremony in disguise, reported to Rome. The quarrel over ``the rites'' lasted a century, one
pope approving and another forbidding, until the mission fell into decline. The conflict between
adaptation and conquest goes on today, Shelley adds, with conquest now usually economic. But this age
produced some of the most creative missionaries of any time.

\begin{keyfacts}
\keypeople{Christopher Columbus; Vasco da Gama; Hernando Cortes; Francisco Pizarro; Charles V;
Bartholomew de Las Casas; Sep\'ulveda; Francis Xavier; Alessandro Valignani; Matthew Ricci; Adam
Schall.}
\keydates{%
  \kd{1492}{Columbus lands in the West Indies (October 12)}
  \kd{1521}{Cortes destroys the Aztec empire}
  \kd{1533}{Pizarro overthrows the Inca empire}
  \kd{1542}{Francis Xavier lands at Goa; Charles V issues the Laws of the Indies}
  \kd{1549}{Xavier lands in Japan}
  \kd{1550}{Las Casas debates Sep\'ulveda at Valladolid}
  \kd{1552}{Xavier dies off the coast of China}
  \kd{1583}{Matthew Ricci allowed to settle in China}
  \kd{1600}{Ricci enters Peking}
  \kd{1614--1642}{Persecution in Japan: 4,045 martyrs}
  \kd{1692}{Chinese imperial edict of toleration for Christianity}}
\keyterms{Age of discovery; adaptation and conquest; \emph{conquistadores}; \emph{encomienda};
Christian prince; Laws of the Indies; mass baptism; Paravars; Nagasaki; rites controversy.}
\end{keyfacts}
